% META: {"id":"1SPE-EXPO-ME-005","chapitre":"1SPE-EXPONENTIELLE","type_objet":"methode","capacites_codes":["C5"],"status":"approved","origin":"generated","status_history":["generated","audited","accepted","approved"],"release_acceptance":"RELEASE_OWNER_BATCH_ACCEPTANCE","acceptance_closure_digest":"sha256:9b3ccf9a81c5520fbb7e03b7b2d3e7bf2057a02834b6d908bf3be5f49f3b553f","acceptance_receipt_digest":"sha256:4fad86f0282e0fa4e0419cca1ad6379ae7af5a280c04a4a90880f90a1c044242"}
\begin{fichemethode}{M5}{Reconnaître les courbes $t\mapsto\mathrm{e}^{kt}$ et $t\mapsto\mathrm{e}^{-kt}$}
\quandUtiliser{J'utilise cette méthode pour associer une courbe à un modèle exponentiel avec $k>0$.}
\pasApas{\begin{enumerate}
  \item Je repère le point commun $(0\,;\,1)$.
  \item Si l'exposant est $kt$, la courbe est croissante.
  \item Si l'exposant est $-kt$, la courbe est décroissante.
  \item Je vérifie que toutes les ordonnées sont strictement positives.
\end{enumerate}}
\exempleRedige{Pour $k=0{,}3$, la courbe de
$t\mapsto\mathrm{e}^{-0{,}3t}$ passe par $(0\,;\,1)$ et décroît lorsque
$t$ augmente.}
\pieges{\begin{itemize}
  \item Faire passer la courbe par l'origine au lieu de $(0\,;\,1)$.
  \item Croire qu'une exponentielle devient négative.
  \item Oublier que $k$ est strictement positif dans cette comparaison.
\end{itemize}}
\verifier{Je calcule la valeur en $0$ et je lis le signe du coefficient de $t$ dans l'exposant.}
\sEntrainer{\refExos{M5}}
\end{fichemethode}
